
\clearpage
\chapter{מבוא לאבטחת שרשרת האספקה בעידן הסוכנים}

\section{עלייתן של מערכות סוכנים אוטומטיות}

מערכות בינה מלאכותית מבוססות סוכנים (\textenglish{LLM Agents}) הפכו לפלטפורמת הוצאה לפועל מרכזית בתחומי אוטומציה עסקית, כתיבת קוד ועיבוד מידע. מסגרות כמו \textenglish{CrewAI} \cite{hong2023metagpt} ו-\textenglish{AutoGen} \cite{wu2023autogen} מאפשרות הרכבה דינמית של כלים (\textenglish{tools/skills}) שסוכן-\textenglish{LLM} טוען בזמן ריצה. גמישות זו מייצרת שטח התקפה חדש: שרשרת האספקה של הכישורים (\textenglish{skill supply chain}).

בדומה לאיומי \textenglish{SolarWinds} ו-\textenglish{Log4Shell} בעולם תוכנה מסורתי, תוקף יכול להחדיר כישורים מורעלים לתוך ריפוזיטורי פומבי, ולגרום לסוכן לטעון קוד זדוני בזמן הביצוע. ההשפעה עלולה לכלול חשיפת סודות (\textenglish{secret exfiltration}), הזרקת פרומפטים (\textenglish{prompt injection}) \cite{perez2022ignore} ושינוי פלטים של הסוכן.

\section{הגדרת האיום}

מודל ה-\textenglish{ClawHavoc} שפיתחנו מתאר תוקף בעל יכולות בינוניות המסוגל:
\begin{itemize}
\item לפרסם כישורים מורעלים לריפוזיטורי \textenglish{PyPI} פומבי
\item להשתמש בשמות דומים לכישורים לגיטימיים (\textenglish{typosquatting})
\item להחדיר הוראות נסתרות בתיאורי הכישור (\textenglish{skill metadata poisoning})
\end{itemize}

מסגרת ה-\textenglish{SkillSieve} שפיתחנו \cite{nakash2024skillsieve} משמשת כשכבת הגנה הבודקת כל כישור לפני טעינתו לצינור הסוכן, ומדגמנת את ספי הסינון האופטימליים.
